% !PRINTED 605
\weberpage{605}
Conversely, in this case 1) the $p$ rational numbers $0, 1, 2, \ldots, p-1$ form a complete system of residues with respect to the modulus $\pi$, and consequently every number in $\mathfrak{o}$, and hence also $\Theta$, is congruent with a rational number. Thus we have

1) \quad $N_a(\pi) = p$ \quad $a_{2,2} = 1$

2) \quad $N_a(\pi) = p^2$ \quad $a_{2,2} = p.$

In the first case $a_{1,2}$ is to be determined from the condition

\weq{3}{a_{1,2} + \Theta \equiv 0 \pmod{\pi}}
while in the second $a_{1,2} = 0$ may be taken. The basis form of $\pi$ thus results, if $u, v$ are the variables, as

1) \quad $\pi = pu + (a_{1,2}+\Theta)v$

2) \quad $\pi = p(u+\Theta v),$

from which one sees that $\pi$ in case 2), as was to be expected, is a \gesperrt{principal functional}.

In order to determine the number $a_{1,2}$ in case 1), we form the norm of $a_{1,2}+\Theta$, which must be divisible by $p$. If $\Theta$ passes over into $\Theta'$ through a change of the sign of $\sqrt{d}$, then by \S.\,164, (8) $\Theta - \Theta' = \sqrt{\Delta}$, and for the norm of $a_{1,2}+\Theta$ we obtain by \S.\,164 (6), a), b):

\weq{4}{N(a_{1,2}+\Theta) = (a_{1,2}+\Theta)(a_{1,2}+\Theta')}
\[ {} = a_{1,2}^2 - d, \quad d\equiv 2, 3\ (\mathrm{mod}\ 4) \]
\[ {} = a_{1,2}^2 + a_{1,2} + \frac{1-d}{4}, \quad d\equiv 1\ (\mathrm{mod}\ 4). \]

If this number is divisible by $p$, then one of the two factors $a_{1,2}+\Theta$, $a_{1,2}+\Theta'$ is divisible by $\pi$. But since $\Theta' = -\Theta$ or $=1-\Theta$, and therefore

a) \quad $a_{1,2} + \Theta' = -(-a_{1,2}+\Theta)$

b) \quad $a_{1,2} + \Theta' = -(-a_{1,2}-1+\Theta)$

one can always take $a_{1,2}$ in such a way that precisely $a_{1,2}+\Theta$ becomes divisible by $\pi$, and $\Theta$ is thus actually congruent with a rational number. The necessary and sufficient condition for $p$ to contain two prime factors is consequently that $a_{1,2}^2 - d$ or $a_{1,2}^2 + a_{1,2} + \frac14(1-d)$ becomes divisible by $p$ by a suitable assumption about $a_{1,2}$. By multiplication by 4 these two conditions pass over into
\[ 4a_{1,2}^2 - 4d \equiv 0, \quad (2a_{1,2}+1)^2 - d \equiv 0 \pmod{4p}, \]
%CONT>>
